% C-written, B-owned labels. Prospective fragment, not included at creation.
\paragraph{A finite Greenberg comparison and its exact infinite-level limit.}
We now give a statement directly about the finite-partition construction,
with its hypotheses separated from the wider historical sheaf programme.
Throughout the following proposition $k$ is algebraically closed,
$R_m=k[t]/(t^m)$, and $m\geq1$.  For a separated reduced finite-type
$k$-scheme $Y$, let $P$ run over its finite partitions into affine locally
closed reduced subschemes and write
\[
 A_P=\prod_{S\in P}\Gamma(S,\mathcal O_S).
\]
Thus $Y_P=\operatorname{Spec}A_P=\coprod_{S\in P}S$.

\begin{proposition}[Proved under stated hypotheses: finite Greenberg comparison]\label{B:finite-Greenberg-piecewise}
Let $X$ be a separated finite-type $R_m$-scheme, not necessarily reduced.
Its finite Greenberg realization $\operatorname{Gr}_m(X)$ is a separated
finite-type $k$-scheme characterized by
\[
 \operatorname{Gr}_m(X)(A)=X(R_m\otimes_k A)
\]
for every $k$-algebra $A$.  Define the finite-partition presheaf by
\[
 N_m(X)(Y)=\varinjlim_P\prod_{S\in P}
 X\bigl(R_m\otimes_k\Gamma(S,\mathcal O_S)\bigr).
\]
There are natural bijections
\[
\begin{aligned}
 N_m(X)(Y)
 &\simeq\operatorname{Hom}_k\bigl(T(Y),\operatorname{Gr}_m(X)\bigr)\\
 &\simeq\operatorname{Hom}_{\operatorname{Abs}_k}
       \bigl(T(Y),T(\operatorname{Gr}_m(X))\bigr).
\end{aligned}
\]
In particular this is a realization of the finite-partition construction
in the absolutely flat category, using its internal products.
The Greenberg functor preserves products and group objects.
\end{proposition}

\begin{proof}
We first supply the elementary equal-characteristic construction.
For an affine $X$, its coordinate ring has a finite presentation
\[
 R_m[z_1,\ldots,z_d]/(f_1,\ldots,f_r).
\]
Replace every $z_i$ by $\sum_{0\leq j<m}z_{ij}t^j$ and expand every
$f_\ell$ modulo $t^m$.  The $m$ coefficient equations for each
$f_\ell$ are polynomial equations over $k$.  Their common zero scheme
represents the required functor on $k$-algebras.  The universal property
makes the construction independent of the chosen presentation.

It glues for a separated finite-type $X$.  Indeed, for every $k$-algebra
$A$, the nilpotent immersion
$\operatorname{Spec}A\to\operatorname{Spec}A[t]/(t^m)$ is a
homeomorphism on underlying spaces.  Consequently a morphism into $X$
factors through an open $U\subset X$ exactly when its reduction does.
Thus the realization of $U$ is the inverse image of $U_k$ under the
reduction morphism.  Applying this to a finite affine cover and its
intersections glues the affine realizations, with precisely the required
functor of points.  The same coefficient construction preserves closed
immersions.  Since the construction preserves products by its universal
property, the realization of the diagonal is closed, proving separatedness.
Finite type follows from the finite affine cover and the explicit finite
polynomial presentations.

This is the equal-characteristic Greenberg functor, equivalently Weil
restriction along $\operatorname{Spec}R_m\to\operatorname{Spec}k$.
The general finite-artinian construction and this identification are
recorded in Proposition--Definition~7.10 and Remarks~7.17(b)--(d) of the
corrected version of Bertapelle--Gonz\'alez-Avil\'es
\cite[pp.~50--54, especially (7.19)]{BBertapelleGonzalezAviles2022}.
Its adjunction is
\[
 \operatorname{Hom}_k(Z,\operatorname{Gr}_m(X))
 \simeq\operatorname{Hom}_{R_m}(Z\times_k\operatorname{Spec}R_m,X).
\]
This proves preservation of products and transports group operations
and their identities.

We have already proved $T(Y)=\varprojlim_PY_P$.  The target
$\operatorname{Gr}_m(X)$ is of finite presentation over $k$, so
finite-presentation descent and equality after refinement give
\[
 \operatorname{Hom}_k(T(Y),\operatorname{Gr}_m(X))
 =\varinjlim_P\operatorname{Hom}_k(Y_P,\operatorname{Gr}_m(X))
\]
by \cite[Lemma~32.10.1(2)--(3), tag 01ZM]{BStacks}.
The adjunction and the finite disjoint union defining $Y_P$ identify
this with $N_m(X)(Y)$.  The second bijection is the defining right
adjunction for $T$.  Although $\operatorname{Gr}_m(X)$ can be
nonreduced, every map from $T(Y)$ or $Y_P$ factors through its
reduction.  No reducedness of the Greenberg realization has been
silently assumed.
\end{proof}

\begin{proposition}[Refuted: interchange of uniform finite partitions with infinite levels]\label{B:uniform-partition-obstruction}
Put $R=k[[t]]$, $A=k[x]$, and let $P$ range over the same finite reduced
affine locally closed partitions of $\operatorname{Spec}A$.  The
natural homomorphism
\[
 D:=\varinjlim_P A_P[[t]]\longrightarrow T(A)[[t]]
\]
is injective but not surjective.  Nevertheless, for every $m\geq1$,
\[
 D/t^mD\simeq T(A)[t]/(t^m),\qquad
 \widehat D^{\,(t)}\simeq T(A)[[t]].
\]
In particular $D$ is a proper dense subring of its $t$-adic completion,
and
\[
 \varinjlim_P\varprojlim_m A_P[t]/(t^m)
 \;\longrightarrow\;
 \varprojlim_m\varinjlim_P A_P[t]/(t^m)
\]
is not surjective.  This is already the additive-group construction for
$\mathbb G_{a,R}$; it is not caused by singularities or non-affineness.
\end{proposition}

\begin{proof}
Each refinement homomorphism $A_P\to A_Q$ is injective: a function
vanishing on every piece of a finite partition vanishes at every prime,
and the original stratum is reduced.  Hence $A_P\to T(A)$ is
injective.  Taking coefficients proves that
$A_P[[t]]\to T(A)[[t]]$ is injective, and the asserted map from $D$
is injective as well.

Choose pairwise distinct $a_0,a_1,\ldots\in k$ and let
\[
 e_n=1-(x-a_n)(x-a_n)^{[-1]}\in T(A).
\]
Here the bracket denotes the unique pointwise inverse.  Thus $e_n$ is
the idempotent equal to $1$ at the closed point $a_n$ and $0$ at all
other primes.  Consider
\[
 z=\sum_{n\geq0}e_nt^n\in T(A)[[t]].
\]
If $z$ came from $A_P[[t]]$, the stratum of $P$ containing the generic
point of the affine line would be a dense open $U$.  Its complement
is finite, so $a_n\in U$ for some $n$.  The $n$th coefficient of the
representative would then be a regular function on $U$ equal to $e_n$.
It vanishes on the dense open $U\setminus\{a_n\}$ and therefore
vanishes identically, since $U$ is reduced.  This contradicts its value
$1$ at $a_n$.  Consequently $z$ is outside the image.

Every finite set of elements of $T(A)=\varinjlim_PA_P$ comes from
one common partition.  Thus $D\to T(A)[t]/(t^m)$ is surjective.
If a representative in $A_P[[t]]$ maps to zero modulo $t^m$, its first
$m$ coefficients are already zero in $A_P$ by injectivity; it is
therefore $t^m$ times another series in $A_P[[t]]$.  The kernel is
exactly $t^mD$.  The asserted quotient and completion formulas follow.
For the specific series $z$, every finite truncation is represented
by the partition isolating the finitely many points involved.  What
fails is the existence of one finite partition for all coefficients.
\end{proof}

\paragraph{Presheaves, sheafification, and the historical scope.}
The last proposition is an exact statement about the displayed
finite-partition presheaves.  It must not be silently promoted to a
statement about a different site or its sheafification.  Such a
promotion requires checking which covers are allowed and whether a
cover can make the series admit one uniform finite partition.  In
particular it does not identify the entire historical target
$\mathcal C'(k)$ or resolve a crossed-out assertion about injective
sheaves.  The finite theorem and the obstruction concern precisely
the stated equal-characteristic, finite-type test setting.

The issue is consistent with, but distinct from, the infinite-level
localization warning in the modern source:
\cite[Proposition~18.2 and Remark~18.3, pp.~99--100]{BBertapelleGonzalezAviles2022}
corrects an earlier proof which had used a finite-level localization
identity at infinite Witt level.  The infinite point formula there
remains true by a different algebraization theorem.  We use neither
that faulty proof nor an unproved interchange of limits.

\begin{proposition}[Refuted: surjectivity after finite-surjective sheafification]\label{B:finite-cover-partition-obstruction}
Let $k$ be algebraically closed.  On the category of reduced affine
finite-type $k$-schemes, equip fibre products with their reduced scheme
structures and take finite surjective morphisms as covers.  For
$Y=\operatorname{Spec}B$ define the presheaves of additive abelian groups
\[
 D(Y)=\varinjlim_P B_P[[t]],\qquad E(Y)=T(B)[[t]],
\]
where $P$ runs over finite affine locally closed reduced partitions of
$Y$.  The induced morphism of associated \emph{sheaves of abelian groups}
\[
 D^\#\longrightarrow E^\#
\]
is not an epimorphism.  The same is true for the associated sheaves of
sets.  In particular the series $z$ of
Proposition~\ref{B:uniform-partition-obstruction} has no local lift from
$D^\#$ on this specified site.  No conclusion about categorical
epimorphisms of sheaves of rings is asserted.
\end{proposition}

\begin{proof}
Finite surjective maps compose and remain finite surjective under base
change and reduction, so these maps form a pretopology.  A finite
jointly surjective family can be replaced by its disjoint union.
For a morphism between affine test schemes, inverse images of affine
locally closed strata are affine locally closed; taking reductions
therefore defines the indicated pullback homomorphisms and presheaves.
The refinement maps are injective by the reducedness argument already
given, so $D\to E$ is injective.

Moreover $E$ is separated for this topology.  If
$\operatorname{Spec}B\to\operatorname{Spec}A$ is finite surjective,
then $T(A)\to T(B)$ is injective.  For each prime $\mathfrak p$ of
$A$ choose a prime $\mathfrak q$ of $B$ above it.  The map
$\kappa(\mathfrak p)\to\kappa(\mathfrak q)$ is injective, and these
are also the corresponding residue fields of $T(A)$ and $T(B)$.
An element in the kernel thus vanishes in every residue field of the
reduced ring $T(A)$, hence is zero.  This argument applies coefficient
by coefficient to $E$, proving separatedness.  Consequently
$E\to E^\#$ is injective.  The same holds for $D$, since $D\subset E$.

We claim that for every finite surjective
$f:Y\to\mathbb A^1_k$, the pullback $f^*z$ lies outside $D(Y)$.
Choose a reduced irreducible component $C\subset Y$ dominating
$\mathbb A^1_k$.  It exists by surjectivity and finiteness.  The
restriction $C\to\mathbb A^1_k$ is finite and dominant, hence
surjective, and $C$ is an integral affine curve.

Suppose that $f^*z$ is represented on a finite partition $P$ of $Y$.
Let $S\in P$ contain the generic point of $C$.  The reduced intersection
$U=C\cap S$ is a dense open in $C$.  Since $C$ is a noetherian curve,
$C\setminus U$ is a finite set of closed points, and its image under
$f$ is finite.  Choose $a_n$ outside that image and a point $y\in C$
over $a_n$; then $y\in U$.  The coefficient of $t^n$ on the stratum
$S$, restricted to $U$, is regular.  Its value at the generic point of
$C$ is zero, since that is the value of $f^*e_n$.  It therefore
vanishes identically on the integral scheme $U$.  At $y$, however,
its residue equals $1$, a contradiction.  This proves the claim.

For completeness we pass from this pointwise claim to the precise
sheaf-theoretic conclusion.  A section of an associated sheaf locally
comes from the presheaf; this follows from the plus construction,
\cite[Lemma~7.10.8 and Theorem~7.10.10]{BStacks}.  If the image of $z$
in $E^\#(\mathbb A^1_k)$ had a local lift to $D^\#$, refining that
cover and composing covers would produce one finite surjective
$f:Y\to\mathbb A^1_k$ and $d\in D(Y)$ whose image is $f^*z$ in
$E^\#(Y)$.  Injectivity of $E(Y)\to E^\#(Y)$ would make the equality
hold already in $E(Y)$, contradicting the claim.  Thus the morphism
is not locally surjective.  This is equivalent to failure of
surjectivity for sheaves of sets, by
\cite[Definition~7.11.1 and Lemma~7.11.2]{BStacks}, and to failure of
epimorphy for sheaves of abelian groups.  The latter can also be seen
from the nonzero associated-sheaf cokernel: the class of $z$ cannot
vanish locally, precisely by the same argument.
\end{proof}

\paragraph{Scope after sheafification.}
This last proposition supplies the additional cover analysis for the
explicit ordinary finite-type site just defined.  It does not assert
that this is every test object or covering allowed in the historical
$\mathcal C'(k)$, or that Grothendieck asserted the failed interchange.
The original phrase specifying finite surjective descent is relevant
motivation, while the test-object adjectives and any larger
pro-affine or perfect setting must be checked separately.
